\section{A bounded cohomological lifting theorem}
\label{sec:bounded}

The quantitative estimate will use bounded cochains on a
transitive $G$-set. We give a norm-controlled lifting argument.
The construction is related to resolutions and relative bounded
cohomology; compare the bounded-cohomological framework of
\cite{Gromov1982}. The treatment of equivariant cochain complexes and boundedly acyclic
actions in \cite[Corollary~5.15 and Remark~5.17]{LiLohMoraschini2026}
is closely related to our construction. Unlike a globally amenable
action, the present coset action has a nonamenable vertex stabilizer;
our argument retains this single-column contribution and gives an
explicit cochain-level bound. We make no claim that the general
lifting mechanism is new.

We introduce the homogeneous and coset cochain complexes.

Let $G$ be a discrete group, $H\leq G$, and $S=G/H$.
The homogeneous complex of bounded real cochains on $G$ is
\[
 C_b^q(G)=\ell^\infty(G^{q+1};\R)^G.
\]
We also consider
\[
 A^p=\ell^\infty_{\mathrm{alt}}(S^{p+1};\R)^G
\]
with the usual alternating simplicial coboundary $\delta$.
Pullback along $\pi:G\to S$ induces a cochain map
$A^\bullet\to C_b^\bullet(G)$, hence
\[
 \rho^k:H^k(A^\bullet)\longrightarrow H_b^k(G;\R).
\]

For all $p,q\geq0$, define
\[
 D^{p,q}=
 \ell^\infty_{\mathrm{alt},S}
 (S^{p+1}\times G^{q+1};\R)^G,
\]
where alternation is imposed only on the $S$-variables.
The horizontal and vertical coboundaries are $\delta$ and $d_G$,
respectively; the total differential on $D^{p,q}$ is
$\delta+(-1)^p d_G$.

\begin{lemma}\label{lem:AW-comparison}
For $p,q\geq0$, let
\[
 R_{p,q}:D^{p,q}\longrightarrow C_b^{p+q}(G)
\]
be the bounded linear map defined by
\[
 (R_{p,q}f)(g_0,\ldots,g_{p+q})
 =f(g_0H,\ldots,g_pH;g_p,\ldots,g_{p+q}).
\]
The maps $R_{p,q}$ have operator norm at most one,
and their sum $R$ on the total complex satisfies
\[
 d_G R=R\bigl(\delta+(-1)^p d_G\bigr).
\]
For $f\in D^{k,0}$ satisfying $d_Gf=0$,
$R f$ is precisely the coset pullback of the
$G$-invariant alternating function determined by $f$.
\end{lemma}

\begin{proof}
Boundedness and $G$-invariance are immediate.
To check the cochain identity, evaluate both sides on
$(g_0,\ldots,g_{p+q+1})$. The horizontal differential
contributes
\[
 \sum_{i=0}^{p+1}(-1)^i
 f(g_0H,\ldots,\widehat{g_iH},\ldots,g_{p+1}H;
   g_{p+1},\ldots,g_{p+q+1}),
\]
and the signed vertical differential contributes
\[
 (-1)^p\sum_{j=0}^{q+1}(-1)^j
 f(g_0H,\ldots,g_pH;
   g_p,\ldots,\widehat{g_{p+j}},\ldots,g_{p+q+1}).
\]
The horizontal term $i=p+1$ cancels the vertical
term $j=0$. All remaining terms are exactly
the homogeneous coboundary of $R_{p,q}f$:
the first $p+1$ delete vertices in the
$S$-variables and the rest delete group variables.
This proves the identity.

If $d_Gf=0$ in vertical degree zero,
$f(s_0,\ldots,s_k;g)$ is independent of $g$,
so $R_{k,0}f$ is the pullback of the corresponding
invariant alternating function of the coset vertices.
\end{proof}

\begin{lemma}\label{lem:horizontal-contraction}
For every $q\geq0$ the horizontally augmented
complex
\[
 0\longrightarrow C_b^q(G)
 \xrightarrow{\jmath}D^{0,q}
 \xrightarrow{\delta}D^{1,q}
 \xrightarrow{\delta}D^{2,q}\longrightarrow\cdots
\]
is contractible by maps of operator norm at most one.
Consequently the augmentation
\[
 \jmath:C_b^\bullet(G)\longrightarrow
          \operatorname{Tot}(D^{\bullet,\bullet}),
 \qquad (\jmath\alpha)(s;g_0,\ldots,g_q)
          =\alpha(g_0,\ldots,g_q)
\]
is a quasi-isomorphism. Under this identification,
the cochain map $R$ of Lemma~\ref{lem:AW-comparison}
is its inverse on cohomology.
\end{lemma}

\begin{proof}
For $p\geq1$ define
\[
 (h^{p,q}f)(s_0,\ldots,s_{p-1};
                         g_0,\ldots,g_q)
 =
 f(g_0H,s_0,\ldots,s_{p-1};g_0,\ldots,g_q).
\]
For $p=0$ use the same formula with no $s$-entries:
\[
 (h^{0,q}f)(g_0,\ldots,g_q)
 =f(g_0H;g_0,\ldots,g_q).
\]
The maps are $G$-equivariant and preserve alternation
in the remaining $S$-variables; their norms are at
most one. Direct expansion of the alternating
coboundary gives
\[
 h^{p+1,q}\delta+\delta h^{p,q}
 =\mathrm{id}_{D^{p,q}}\quad(p\geq1),
 \qquad
 h^{1,q}\delta+\jmath h^{0,q}
 =\mathrm{id}_{D^{0,q}},
 \qquad h^{0,q}\jmath=\mathrm{id}.
\]
The formulas involve $g_0H$, which is why
the homotopies are equivariant despite the
absence of a $G$-fixed vertex of $S$.

Thus the horizontal augmentation is contractible
in every fixed vertical degree. Filtering the
first-quadrant total complex by the vertical
degree, its horizontal-first spectral sequence
has only the augmented-column contribution
$C_b^\bullet(G)$; equivalently, the usual
bounded double-complex comparison lemma
applies to this explicitly contractible
horizontal resolution. This proves that $\jmath$
is a quasi-isomorphism. Finally
$R\jmath=\mathrm{id}_{C_b^\bullet(G)}$
on cochains, so $R$ induces the inverse.
\end{proof}

\begin{theorem}\label{thm:bounded-lift}
Assume that the pointwise stabilizer of any two distinct
elements of $S$ is amenable. Then for each $k\geq2$
\begin{equation}\label{eq:bounded-image}
 \operatorname{im}\rho^k
 =\ker\bigl(H_b^k(G;\R)\longrightarrow H_b^k(H;\R)\bigr).
\end{equation}
More quantitatively, if $\alpha\in C_b^k(G)$ is a
homogeneous cocycle satisfying
$\alpha|_{H^{k+1}}=0$, there is a cocycle $\beta\in A^k$
such that
\[
 [\pi^*\beta]=[\alpha],\qquad
 \|\beta\|_\infty\leq k(k+1)!\|\alpha\|_\infty.
\]
\end{theorem}

We first record the bounded contractions used in the proof.

\begin{lemma}\label{lem:bounded-vertical}
For $p\geq1$, the vertical homogeneous coboundary complex
$D^{p,\bullet}$ has bounded contracting homotopies
in positive degrees, of norm at most one.
\end{lemma}

\begin{proof}
Consider an ordered tuple of pairwise distinct vertices
$(s_0,\ldots,s_p)$. Its pointwise stabilizer $K$ fixes
at least two distinct vertices, so is amenable.
The setwise stabilizer is a finite extension of $K$
and therefore also amenable. On each orbit of such
tuples, equivariant cochains are identified with bounded
cochains on $G^{q+1}$ invariant under the tuple stabilizer,
with the sign representation accounting for odd vertex
permutations.

We make the contraction explicit. Let $K_\sigma$
be the setwise stabilizer of the unordered vertex set
of an ordered tuple $\sigma$, and let
$\chi_\sigma:K_\sigma\to\{+1,-1\}$ be its permutation-sign
character. Fix a left $K_\sigma$-invariant mean
$m_\sigma$ on $\ell^\infty(G)$. Such a mean exists:
$K_\sigma$ is amenable, and one may choose a mean on
each $K_\sigma$-orbit of $G$ and then a mean on the
set of these orbits. On the vertical complex associated
with $\sigma$, cochains satisfy
\[
 f(kg_0,\ldots,kg_q)
  =\chi_\sigma(k)f(g_0,\ldots,g_q).
\]
Define, for $q\geq1$,
\[
 (h_\sigma f)(g_0,\ldots,g_{q-1})
  =m_\sigma\bigl(x\mapsto
       f(x,g_0,\ldots,g_{q-1})\bigr).
\]
Left invariance of $m_\sigma$ shows that
$h_\sigma f$ obeys the same $\chi_\sigma$-equivariance
as $f$. The standard insertion identity
\[
 d_Gh_\sigma+h_\sigma d_G=\mathrm{id}
\]
holds in positive vertical degrees, and
$\|h_\sigma\|\leq1$. The sign character is
handled by equivariance; no additional
averaging factor is required.

Choose an orbit representative for each unordered
tuple of distinct vertices and transport this
construction equivariantly. The resulting
homotopies have uniform norm at most one.
Alternation forces cochains on tuples with
repeated vertices to vanish.
\end{proof}

\begin{lemma}\label{lem:relative-primitive}
Let $\alpha\in C_b^k(G)$, $k\geq1$, be a bounded
homogeneous cocycle with zero restriction as a cochain
to $H^{k+1}$. The corresponding element of
$D^{0,k}$ admits a vertical primitive $b_0\in
D^{0,k-1}$ satisfying
\[
 d_Gb_0=\alpha,\qquad
 \|b_0\|_\infty\leq k\|\alpha\|_\infty.
\]
\end{lemma}

\begin{proof}
Choose an $H$-equivariant retraction of sets
$r:G\to H$ that is the identity on $H$; choose
representatives of the left $H$-orbits, with the
representative $1$ for $H$ itself. The standard
prism homotopy between the identity and $r$ on
the homogeneous bar complex has, on cochains
of degree $k$, at most $k$ summands. Consequently
its operator norm is at most $k$.
Since $\alpha$ restricts to zero on $H$,
the prism identity $dK+Kd=\mathrm{id}-r^*$
gives $d_G(K\alpha)=\alpha$.
Transporting this primitive from the vertex $H$
to the other vertices of $G/H$ gives the
required $b_0$.
\end{proof}

\begin{proof}[Proof of Theorem~\ref{thm:bounded-lift}]
We use the first-quadrant double complex
$D^{p,q}$ with horizontal differential
$\delta$ and vertical differential $d_G$, and
total differential $d_{\mathrm{tot}}=
\delta+(-1)^pd_G$ on $D^{p,q}$.
The horizontal contraction in
Lemma~\ref{lem:horizontal-contraction}
shows that $\jmath:C_b^\bullet(G)\to
\operatorname{Tot}D$ is a quasi-isomorphism.
The map $R$ of Lemma~\ref{lem:AW-comparison}
is a cochain left inverse to $\jmath$.
In particular, a cocycle in
\[
 A^p=\ker(d_G:D^{p,0}\to D^{p,1})
\]
has image under $R$ equal to its coset
pullback.

First suppose that $\alpha\in C_b^k(G)$
is a cocycle whose restriction to $H^{k+1}$
vanishes as a cochain. Write
$a_0=\jmath\alpha\in D^{0,k}$.
Lemma~\ref{lem:relative-primitive}
provides $b_0\in D^{0,k-1}$ with
\[
 d_Gb_0=a_0,\qquad
 \|b_0\|_\infty\leq k\|\alpha\|_\infty.
\]
Set $a_1=-\delta b_0\in D^{1,k-1}$.
Then $a_0-d_{\mathrm{tot}}b_0=a_1$,
and $\delta a_1=d_Ga_1=0$.

For $1\leq p<k$, suppose inductively
that $a_p\in D^{p,k-p}$ is closed under
both $\delta$ and $d_G$. Its vertical
degree $k-p$ is positive, so
Lemma~\ref{lem:bounded-vertical}
gives a primitive
\[
 b_p=(-1)^ph_pa_p\in D^{p,k-p-1},
 \qquad d_Gb_p=(-1)^pa_p,
 \qquad \|b_p\|_\infty\leq\|a_p\|_\infty.
\]
Put $a_{p+1}=-\delta b_p$. Then
\[
 a_p-d_{\mathrm{tot}}b_p=a_{p+1},
 \qquad
 \delta a_{p+1}=0,
 \qquad
 d_Ga_{p+1}=0.
\]
The last identity follows from
$d_G\delta=\delta d_G$ and
$\delta a_p=0$. By induction
$a_k\in D^{k,0}$ is closed under
both differentials, so it is an element
$\beta\in A^k$ with $\delta\beta=0$.
Summing the displayed total-coboundary
identities and applying $R$ proves
$[\pi^*\beta]=[\alpha]$.

On $D^{p,q}$ the horizontal coboundary
has operator norm at most $p+2$.
Consequently
\[
 \|\beta\|_\infty
 \leq (k+1)!\|b_0\|_\infty
 \leq k(k+1)!\|\alpha\|_\infty.
\]

Now suppose only that
$[\alpha]|_H=0$ in bounded cohomology.
There is a bounded homogeneous
$H$-invariant $(k-1)$-cochain
$\theta$ on $H$ with
$d_G\theta=\alpha|_{H^{k+1}}$.
Define a bounded homogeneous
$G$-invariant cochain $\widehat\theta$
on $G^k$ as follows. If all its
arguments lie in a common left coset
$tH$, set
\[
 \widehat\theta(th_0,\ldots,th_{k-1})
    =\theta(h_0,\ldots,h_{k-1});
\]
otherwise set its value to zero.
This is well-defined: replacing $t$
by $th$ translates all arguments
of $\theta$ simultaneously by $h^{-1}$,
and $\theta$ is $H$-invariant.
The resulting cochain is bounded
and $G$-invariant. Thus
$\alpha'=\alpha-d_G\widehat\theta$
has zero cochain restriction to $H$.
Applying the first part to $\alpha'$
gives a cocycle in $A^k$ mapping to
$[\alpha]$.

Conversely, any alternating cocycle on
$S$ pulls back to a cocycle vanishing
on $H^{k+1}$, because all those
arguments have coset $H$. Hence
$\operatorname{im}\rho^k$ is exactly
the restriction kernel, proving
\eqref{eq:bounded-image}.
\end{proof}

We specialize the lifting theorem to the diagonal subgroup.

For the application let $\Gamma$ be the fundamental group
of a closed negatively curved manifold and put
\[
 G=\Gamma\times\Gamma,\qquad H=\Delta\Gamma.
\]
The $G$-set $G/H$ is naturally isomorphic to $\Gamma$,
with $(g,h)a=gah^{-1}$.

\begin{lemma}\label{lem:diagonal-stabilizers}
The common stabilizer in $G$ of two distinct labels
$a,b\in\Gamma$ is isomorphic to
$C_\Gamma(b^{-1}a)$. Hence it is infinite cyclic.
\end{lemma}

\begin{proof}
An element $(g,h)$ fixes $a$ if and only if
$g=aha^{-1}$; it fixes $b$ as well exactly when
$h$ commutes with $b^{-1}a$.
Since $b^{-1}a\ne1$, its centralizer in the
torsion-free hyperbolic group $\Gamma$ is
infinite cyclic \cite{BridsonHaefliger1999}. This identifies the
common stabilizer.
\end{proof}

\begin{proposition}\label{prop:difference-lift-bound}
Let $\omega\in C_b^n(\Gamma)$ be an alternating homogeneous
$n$-cocycle, $n\geq2$, and put
\[
 D=p_2^*\omega-p_1^*\omega\in C_b^n(G).
\]
There is a bounded alternating $G$-invariant cocycle
$\widetilde D$ on $\Gamma^{n+1}$ with
\[
 [\pi^*\widetilde D]=[D],\qquad
 \|\widetilde D\|_\infty
 \leq n(n+1)!\,\|\omega\|_\infty.
\]
\end{proposition}

\begin{proof}
For the initial vertical contraction at the coset
$H=\Delta\Gamma$, use the $H$-equivariant retraction
\[
 r:G\longrightarrow H,\qquad r(g,h)=(g,g).
\]
The homogeneous prism cochain of degree $n-1$ is
\[
 KD(g_0,\ldots,g_{n-1})
 =\sum_{i=0}^{n-1}(-1)^i
 D(rg_0,\ldots,rg_i,g_i,\ldots,g_{n-1}).
\]
Here the expression in the $i$-th summand has
$n+1$ group arguments. The first-coordinate
parts of $rg_i$ and $g_i$ agree. Since $\omega$
is alternating, the $p_1^*\omega$ contribution
of every summand is zero. Consequently
\[
 \|KD\|_\infty\leq n\|\omega\|_\infty.
\]
The prism identity and $D|_{H^{n+1}}=0$
give $d_G(KD)=D$.

Transport $KD$ equivariantly to the other vertex
cosets, obtaining $b_0\in D^{0,n-1}$
with the same bound. In the double complex
of Theorem~\ref{thm:bounded-lift}, set
\[
 a_1=-\delta b_0,\qquad
 b_p=(-1)^p h_pa_p,\qquad
 a_{p+1}=-\delta b_p
 \quad(1\leq p<n),
\]
where $h_p$ denotes the norm-one vertical
contraction in column $p$. This removes the
positive vertical degrees successively.
The horizontal differential from column
$p$ to column $p+1$ has norm at most $p+2$.
It follows that the final cocycle
$\widetilde D=a_n$ satisfies
\[
 \|\widetilde D\|_\infty
 \leq n(2\cdot3\cdots(n+1))
       \|\omega\|_\infty
 =n(n+1)!\|\omega\|_\infty.
\]
The total-complex coboundary identities imply
$[\pi^*\widetilde D]=[D]$, as in
Theorem~\ref{thm:bounded-lift}.
\end{proof}

\begin{corollary}\label{cor:bounded-diagonal}
For $k\geq2$, every bounded cohomology class on
$\Gamma\times\Gamma$ restricting trivially to
$\Delta\Gamma$ is represented by the pullback
of a bounded alternating invariant cocycle
on $\Gamma^{k+1}$. A cocycle whose restriction
vanishes literally admits the norm estimate
of Theorem~\ref{thm:bounded-lift}.
\end{corollary}

\begin{remark}
The norm estimate concerns a cocycle with zero
cochain restriction. It is not a uniform bound
in terms of the norm of an arbitrary class in
the kernel unless a norm-controlled choice of
relative representative is supplied.
This distinction is important in Section 5,
where the cocycles are chosen to vanish
explicitly on the diagonal.
\end{remark}
